% TODO make hw stylesheet
\documentclass[10pt]{article}
\usepackage[letterpaper,top=1in]{geometry}
\usepackage{quiver}
\usepackage[dvipsnames,svgnames]{xcolor}
\usepackage[shortlabels]{enumitem}
\usepackage[framemethod=TikZ]{mdframed}
\usepackage{amsmath,amssymb,amsthm}
\usepackage[colorlinks]{hyperref}
\usepackage{multicol}
\usepackage{tabularx}
\usepackage{stmaryrd}

\hypersetup{
	colorlinks=true,
	linkcolor=blue,
	filecolor=magenta,
	urlcolor=magenta,
	pdftitle={MAT 534 HW}
}

\usepackage{mathtools}
\usepackage{thmtools}
\usepackage[textsize=footnotesize]{todonotes}
\usepackage{titling}

\reversemarginpar
\mdfdefinestyle{mdbluebox}{roundcorner=10pt,innerbottommargin=9pt,
	linecolor=blue,backgroundcolor=TealBlue!5,}
\declaretheoremstyle[headfont=\sffamily\bfseries\color{MidnightBlue},
mdframed={style=mdbluebox},]{thmbluebox}

\mdfdefinestyle{mdbloobox}{frametitlefont=\bfseries,innerbottommargin=8pt,
	nobreak=true,backgroundcolor=TealBlue!5,linecolor=blue,}
\declaretheoremstyle[headfont=\bfseries\color{MidnightBlue},
mdframed={style=mdbloobox},headpunct={\\[3pt]},postheadspace=0pt,]{thmbloobox}

\mdfdefinestyle{mdredbox}{frametitlefont=\bfseries,innerbottommargin=8pt,
	nobreak=true,backgroundcolor=Salmon!5,linecolor=RawSienna,}
\declaretheoremstyle[headfont=\bfseries\color{RawSienna},
mdframed={style=mdredbox},headpunct={\\[3pt]},postheadspace=0pt,]{thmredbox}

\mdfdefinestyle{mdgreenbox}{linecolor=ForestGreen,backgroundcolor=ForestGreen!5,
	linewidth=2pt,rightline=false,leftline=true,topline=false,bottomline=false,}
\declaretheoremstyle[headfont=\bfseries\sffamily\color{ForestGreen!70!black},
mdframed={style=mdgreenbox},headpunct={ --- },]{thmgreenbox}

\mdfdefinestyle{mdorangebox}{linecolor=RedOrange,backgroundcolor=RedOrange!5,
	linewidth=2pt,rightline=false,leftline=true,topline=false,bottomline=false,}
\declaretheoremstyle[headfont=\bfseries\sffamily\color{RedOrange!70!black},
mdframed={style=mdorangebox},headpunct={ --- },]{thmorangebox}

\mdfdefinestyle{mdblackbox}{linecolor=black,backgroundcolor=RedViolet!5!gray!5,
	linewidth=3pt,rightline=false,leftline=true,topline=false,bottomline=false,}
\declaretheoremstyle[mdframed={style=mdblackbox}]{thmblackbox}

\declaretheoremstyle[spaceabove=6pt,spacebelow=6pt]{basehead}

\declaretheorem[style=basehead,name=Problem]{problem}
\declaretheorem[style=thmredbox,name=Problem,numbered=no]{problem*}
\declaretheorem[style=thmbloobox,name=Setup]{setup} % for config geo
\declaretheorem[style=thmredbox,name=Example,sibling=problem]{example}
\declaretheorem[style=thmredbox,name=$\bigstar$ Problem,sibling=problem]{niceprob}
\declaretheorem[style=thmbluebox,name=Theorem,sibling=problem]{theorem}
\declaretheorem[style=thmbluebox,name=Corollary,sibling=theorem]{corollary}
\declaretheorem[style=thmbluebox,name=Proposition,sibling=theorem]{prop}
\declaretheorem[style=thmbluebox,name=Lemma,numberwithin=problem]{lemma}
\declaretheorem[style=thmbluebox,name=Theorem,numbered=no]{theorem*}
\declaretheorem[style=thmbluebox,name=Lemma,numbered=no]{lemma*}
\declaretheorem[style=thmgreenbox,name=Claim,numberwithin=problem]{claim}
\declaretheorem[style=thmgreenbox,name=Claim,numbered=no]{claim*}
\declaretheorem[style=thmblackbox,name=Remark,numberwithin=problem]{remark}
\declaretheorem[style=thmblackbox,name=Remark,numbered=no]{remark*}
\declaretheorem[style=thmgreenbox,name=Definition,sibling=theorem]{definition}
\declaretheorem[style=thmgreenbox,name=Definition,numbered=no]{definition*}

\providecommand{\ol}{\overline}
\providecommand{\quiv}{\equiv}
\providecommand{\ceil}[1]{\left\lceil #1 \right\rceil}
\providecommand{\floor}[1]{\left\lfloor #1 \right\rfloor}
\newcommand{\dd}{\mathrm{d}}
\providecommand{\ul}{\underline}
\providecommand{\eps}{\varepsilon}
\providecommand{\half}{\frac{1}{2}}
\providecommand{\CC}{\mathbb C}
\providecommand{\FF}{\mathbb F}
\providecommand{\NN}{\mathbb N}
\providecommand{\QQ}{\mathbb Q}
\providecommand{\RR}{\mathbb R}
\providecommand{\ZZ}{\mathbb Z}
\providecommand{\dg}{^\circ}
\providecommand{\ii}{\item}
\providecommand{\setdiff}{\smallsetminus}
\providecommand{\alert}[1]{{\sffamily\textbf{\textcolor{blue}{#1}}}}
\providecommand{\inv}{^{-1}}
\providecommand{\mb}[1]{\mathbf{#1}}

\newenvironment{soln}{\begin{proof}[Solution]}{\end{proof}}

\providecommand{\pow}{\operatorname{Pow}}
\providecommand{\aut}{\operatorname{Aut}}
\providecommand{\vocab}[1]{{\textbf{\textcolor{ForestGreen}{#1}}}}
\providecommand{\lcm}{\operatorname{lcm}}
\providecommand{\abs}[1]{\left|#1\right|}
\providecommand{\Ob}{\operatorname{Ob}}
\providecommand{\Hom}{\operatorname{Hom}}
\providecommand{\Aut}{\operatorname{Aut}}
\providecommand{\id}{\operatorname{id}}
\newcommand{\doubleparen}[1]{(\!(#1)\!)}

\usepackage{mathrsfs}

% place title at top right
\setlength{\droptitle}{-8ex}
\pretitle{\begin{flushright}\Large\bfseries}
\posttitle{\par\end{flushright}}
\preauthor{\begin{flushright}}
\postauthor{\end{flushright}}
\predate{\begin{flushright}}
\postdate{\end{flushright}}

\title{MAT 534 HW07}
\author{\large Aditya Pahuja}
\date{\large Due 11/05}
\begin{document}
\maketitle

\begin{problem}
    Let $\varphi\colon R\to S$ be a homomorphism of unital commutative rings.
    Show that if $P$ is a prime ideal in $S$,
    then $\varphi\inv(P)$ is a prime ideal in $R$.
\end{problem}

\begin{soln}
    Let $\psi\colon S\to S/P$ be the projection homomorphism.
    Then $\ker\psi = P$, so $\ker(\psi\circ\varphi) = \varphi\inv(P)$,
    and of course $\psi\circ\varphi\colon R\to S/P$ is surjective,
    so $S/P\cong R/\varphi\inv(P)$.
    Since $P$ is prime in $S$, $S/P$ is an integral domain,
    hence so is $R/\varphi\inv(P)$, implying $\varphi\inv(P)$ is
    a prime ideal of $R$.
\end{soln}

\begin{problem}
    Let $G$ be a finite subgroup of the group of units $F^\times$ of field $F$.
    \begin{enumerate}[(a)]
        \ii Let $d$ be a divisor of $\abs G$.
	Show that $G$ contains at most $d$ elements whose order divides $d$.
	\ii Let $\abs G = p_1^{e_1}\cdots p_r^{e_r}$
	be the prime factorization of $\abs G$.
	Show that for each $i=1,\dots,r$, there is an element
	of order $p_i^{e_i}$ in $G$.
	\ii Show that $G$ contains an element of order $|G|$,
	and conclude that $G$ is a cyclic group.
    \end{enumerate}
\end{problem}

\begin{soln}
    (a) The polynomial $x^d - 1$ has at most $d$ roots
    because its degree is $d$; this is independent of
    $d$ being a divisor of $\abs G$.
    \\

    (b) Since $G$ is a finite abelian group, it is a direct product
    of abelian groups $G_1\times\cdots\times G_r$ where $\abs{G_i} = p_i^{e_i}$.
    Let $p_i^{k_i}$ be the largest power of $p_i$ that appears as
    the order of some element of $G$; since there are $p_i^{e_i}$ elements
    (corresponding to the elements of $G_i$) whose order is a power of $p_i$,
    we must have $p_i^{k_i}\ge p_i^{e_i}$ by part (a),
    so $k_i \ge e_i$ (hence equals $e_i$), proving the existence
    of an element with order $p_i^{e_i}$.
    \\

    (c) $G$ is a finite abelian group, so it is
    a direct product of cyclic groups; these cyclic groups must be
    the $r$ groups $\ZZ/p_i^{e_i}\ZZ$ by the order information, so
    $G\cong\ZZ/p_1^{e_1}\ZZ\times\cdots\times\ZZ/p_r^{e_r}\ZZ \cong \ZZ/|G|\ZZ$
    by the Chinese remainder theorem.
\end{soln}

\begin{problem}
    Let $\FF_p$ be the field with $p$ elements.
    \begin{enumerate}[(a)]
	\ii Let $g(x)\in\FF_p[x]$ be an irreducible polynomial with degree $n\ge1$.
	Show that $F=\FF_p[x]/(g(x))$ is a field with $q=p^n$ elements.
	\ii Show that $g(x)$, when considered as
	a polynomial in $F[x]$, has a root in $F$.
	\ii Show that every element of $F$ is a root of the polynomial $x^q - x$.
	\ii Show that, in the polynomial ring $\FF_p[x]$,
	the $\gcd$ of the two polynomials $g(x)$ and $x^q - x$ is not $1$.
	\ii Conclude that $g(x)$ divides $x^q - x$ in the ring $\FF_p[x]$.
    \end{enumerate}
\end{problem}

\begin{soln}
    (a) $F$ is a field because $\FF_p[x]$ is a PID
    and irreducible elements of a PID generate maximal ideals.

    Furthermore, if $\alpha$ is the image of $x$
    under the projection map from $\FF_p[x]$ to $F$,
    then $1$, $\alpha$, \dots, $\alpha^{n-1}$ spans
    $F$ as a $\FF_p$-vector space, since any polynomial is in the same
    coset of $(g(x))$ as its remainder upon division by $g(x)$,
    which has degree at most $n-1$.

    Additionally, $1$, $\alpha$, \dots, $\alpha^{n-1}$ are
    linearly independent over $\FF_p$
    because linear dependence would imply that
    $\alpha$ is the root of a polynomial $h$ of degree $n-1$
    with coefficients in $\FF_p$, which means $h(x)\in\FF_p[x]$
    is in $(g(x))$; this is a contradiction.
    \\

    (b) The element $\alpha\in F$ mentioned in (a) is a root of $g$
    because $g(x)$ and $0$ are in the same coset of $(g(x))$, i.e.
    the images $g(\alpha)$ and $0$ are equal inside $F$.
    \\

    (c) The multiplicative group of $F$ is finite with order $q-1$,
    so by the previous problem it is cyclic of order $q-1$.
    Thus $a^{q-1} = 1$ if $a\ne 0$, so
    $a(a^{q-1}-1) = 0$ for all $a\in F$, i.e. $x^q - x$ has
    every $a\in F$ as a root.
    \\

    (d) The projection homomorphism $x\mapsto \alpha$ onto $F$ sends
    $x^q - x$ to $\alpha^q - \alpha = 0$, so $x^q - x$
    lies in the ideal $(g(x))$, hence is divisible by $g(x)$.
    \\

    (e) uhhhh see part (d)?
\end{soln}

\begin{problem}
    Let $F$ be a field. Prove that the additive group $(F,+)$
    is not isomorphic to the multiplicative group $(F^\times,\cdot)$.
\end{problem}

\begin{soln}
    Let $\varphi\colon F^\times\to F$ be a homomorphism.
    Then $\varphi(-1) + \varphi(-1) = \varphi(1) = 0$
    implies $\varphi(-1) = \varphi(1) = 0$,
    so $\varphi$ is not bijective, hence cannot be an isomorphism.
\end{soln}

\begin{problem}
    Show that the polynomial $(x-1)(x-2)\cdots(x-n)+1$
    is irreducible in the ring $\ZZ[x]$ for all $n\ge 1$ except $n=4$.
\end{problem}

\begin{soln}
    Let $f(x) = (x-1)(x-2)\cdots(x-n)+1$ and
    suppose $f(x) = p(x)q(x)$ for some $p(x),q(x)\in\ZZ[x]$
    where $p$ and $q$ both have degree less than $n$.
    Then $p(k)q(k) = 1$ for all $k=1,2,\dots,n$, so
    $p(k) = q(k) = \pm 1$. Since $p$ and $q$ agree
    at $n > \max\{\deg p,\deg q\}$ points, $p$ and $q$ are the same polynomial.

    Therefore $f(x) = p(x)^2$, so if $f$ is reducible
    then $n$ must be even and $d \coloneq \deg p$ is $n/2$.
    This implies $p(k) = 1$ has $n/2$ solutions in $\{1,2,\dots,n\}$
    (otherwise either $p(k) - 1$ or $p(k) + 1$ has more than $n/2$ roots);
    let $R = \{r_1,\dots,r_d\}$ be the set of these solutions
    and $S = \{s_1,\dots,s_d\}$ the remaining elements of $\{1,2,\dots,n\}$
    (i.e. the solutions to $p(k) = -1$).
    Since $f(x)$ is monic, $p(x)$ is monic, so $p(x)$ is uniquely determined as
    \begin{equation*}
        p(x) = (x - r_1)(x - r_2)\cdots(x - r_d) + 1.
    \end{equation*}
    There exists $s\in S$ such that $r - s \ge d$ for some $r\in R$
    because if $1\in R$ then $R$ contains at most $d-1$ of the
    numbers $d+1$, $d+2$, \dots, $d+d$, hence one of them is in $S$
    (and vice versa if $1\in S$). For such an $s$, we conclude that
    \begin{equation*}
	2 = \abs{p(s) - 1} = \abs{(s - r_1)(s - r_2)\cdots(s - r_d)}
	\ge \abs{r-s} \ge d,
    \end{equation*}
    so if $d > 2$ (hence $n > 4$) then $f$ is irreducible.
    If $d = 1$ then $\abs{s - r_1} = 1 \ne 2$,
    so $f$ is still irreducible when $n=2$.
    (When $d = 2$, $f$ is reducible: $(x-1)(x-2)(x-3)(x-4) + 1
    =(x^2 - 5x + 5)^2$.)
\end{soln}

\pagebreak

\begin{problem}
    Let $F$ be a field. Let $R$ be the set of polynomials in $F[x]$
    whose coefficient of $x$ is equal to $0$.
    Show that $R$ is a subring of $F[x]$, and that $R$ is not a UFD.
\end{problem}

\begin{soln}
    To show that $R$ is a subring of $F$ we verify that
    \begin{itemize}
        \ii $1\in R$ trivially;
	\ii if $f(x),g(x)\in R$ then the coefficient of $x$ in
	$f(x) + g(x)$ is $0 + 0 = 0$; and
	\ii if $f(x),g(x)\in R$ then the coefficient of $x$ in
	$f(x)g(x)$ is $c_f\cdot 0 + 0\cdot c_g = 0$
	where $c_f,c_g$ are the constant terms of $f(x)$, $g(x)$ respectively.
    \end{itemize}
    However, $R$ is not a UFD because
    $x^6 = x^2\cdot x^2\cdot x^2 = x^3\cdot x^3$
    has two factorizations into irreducibles;
    both $x^2$, $x^3$ are irreducible because any factorization
    of either one must have a linear factor.
\end{soln}

\begin{problem}
    Show that the ideal $(x,y)^n = (x^n, x^{n-1}y, \dots, y^n)$
    in the ring $\QQ[x,y]$ cannot be generated by fewer than $n+1$ elements.
\end{problem}

\begin{soln}
    For $r\in\ZZ_{\ge0}$, let $I_r$ be the ideal $(x,y)^r$
    and suppose that $I_n$ is generated by the polynomials
    $p_1(x,y),\dots,p_m(x,y)\in\QQ[x,y]$.
    Since $p_i(x,y)\in I_n$, each monomial of $p_i$
    has degree at least $n$.

    The quotient of $\QQ[x,y]$-modules $I_n/I_{n+1}$ is isomorphic
    as a $\QQ$-vector space to the collection of
    homogeneous polynomials in $\QQ[x,y]$
    (since the quotient kills off all monomials with degree at least $n+1$,
    and $I_n$ a priori has all monomials appearing with degree at least $n$),
    so it has dimension $n + 1$ as a $\QQ$-vector space
    (one example of a basis is $x^n$, \dots, $xy^{n-1}$, $y^n$).
    On the other hand, an element $p(x,y)$ of $I_n$ can be written as
    \begin{equation*}
	p(x,y) = \sum_{i=1}^m q_i(x,y)p_i(x,y)
    \end{equation*}
    for $q_1(x,y),\dots,q_m(x,y)\in\QQ[x,y]$.
    When the non-constant terms of $q_i(x,y)$
    multiply with $p_i(x,y)$, they produce monomials that have degree
    at least $n+1$, so, projecting onto $I_n/I_{n+1}$,
    \begin{equation*}
	p(x,y) = \sum_{i=1}^m c_i\cdot p_i(x,y)
    \end{equation*}
    where $c_i$ is the constant term of $q_i(x,y)$.
    Thus each element of $I_n/I_{n+1}$ is manifested as a $\QQ$-linear
    combination of the $p_i$, implying that the $p_i$ span
    the $(n+1)$-dimensional vector space $I_n/I_{n+1}$, so $m\ge n+1$.
\end{soln}

\begin{problem}
    Let $R$ be a UFD, and $F$ its field of fractions.
    \begin{enumerate}[(a)]
        \ii Show that if $p\in R$ is irreducible,
	then the constant polynomial $p$ is irreducible in $R[x]$.
	\ii Show that if $f(x)\in R[x]$ is primitive and irreducible in $F[x]$,
	then $f(x)$ is irreducible in $R[x]$.
	\ii Suppose that we have two factorizations
	\begin{equation*}
	    p_1\cdots p_k\cdot f_1(x)\cdots f_m(x) =
	    q_1\cdots q_\ell\cdot g_1(x)\cdots g_n(x)
	\end{equation*}
	in the ring $R[x]$, with $p_i,q_i\in R$ irreducible
	and $f_i(x),g_i(x)\in R[x]$ primitive and irreducible in $F[x]$.
	Prove that $k=\ell$, $m=n$,
	and that the factors on both sides are equal up to reordering
	and multiplication by a unit in $R$.
    \end{enumerate}
\end{problem}

\begin{soln}
    (a) If $p$ has a factorization into non-units in $R[x]$,
    the factors both have degree zero, so the factors would be in $R$;
    $p$ being irreducible in $R$ means such a factorization doesn't exist.
    \\

    (b) If $f(x)$ has a factorization into non-units,
    this factorization must be of the form $p(x)q(x)$
    for some non-constant $p(x),q(x)\in R[x]$
    (we can assume non-constant because $f$ is primitive).
    However, this factorization is then also a factorization
    into non-units in $F[x]$, contradicting the irreducibility of $f(x)$.
    \\

    (c) By multiplying each polynomial by a unit
    (and absorbing the unit into one of the $p_i$, $q_i$)
    we can force all the $f_i$ and $g_i$ to be monic,
    so $p_1\cdots p_k = q_1\cdots q_\ell$
    and $f_1(x)\cdots f_m(x) = g_1(x)\cdots g_n(x)$.
    The former equation, together with $R$ being a UFD,
    implies that the $p_i$ and $q_i$ are the same up to reordering.

    For the $f_i(x)$ and $g_i(x)$, we induct on $m$:
    specifically, we show that if an element of $R[x]$
    is expressed as the product of $m$ irreducible primitive polynomials,
    then this product is unique. This is trivially true at $m=1$,
    so assume that the conclusion holds for all $k\le m$.
    Since $F[x]$ is a PID, irreducible elements are prime, so
    $f_m(x)\mid g_1(x)\cdots g_n(x)$ implies that
    $f_m(x)$ divides, hence is equal (up to a unit in $F$) to one of the $g_i$,
    say WLOG $g_n(x)$ by relabeling. Since the $f_i$, $g_i$
    are all primitive, this unit in $F$ is actually a unit in $R$,
    so we can cancel $f_n$ and $g_m$
    to get that $f_1(x)\cdots f_{m-1}(x) = g_1(x)\cdots g_{n-1}(x)$,
    at which point the inductive hypothesis tells us that the multisets
    $\{f_1,\dots,f_{m-1}\}$ and $\{g_1,\dots,g_{n-1}\}$ are equal,
    finishing the proof.
\end{soln}










\end{document}
